\begingroup
\renewcommand{\arraystretch}{1.45}
\fancyhead[L]{\small TRANSZENDENTE FUNKTIONEN}
\Needspace{43\baselineskip}
\AGMArticle{10}
Die folgenden Beispiele werden zeigen, wie schnell diese Reihen konvergieren:

\Needspace{13\baselineskip}\par I. Es sei \(x=\surd 2,\ y=1\) (vgl. Artikel 3). Hier gilt \(xx-yy=1\)
\begingroup\small\setlength{\arraycolsep}{3pt}
\[
\begin{array}{rcll}
2(x'x'-y'y')&=\tfrac12(x-y)^2=&0{,}0857864376&2690\\[7pt]
4(x''x''-y''y'')&=(x'-y')^2=&0{,}\mathord{.}\mathord{.}\mathord{.}3203980&4927\\[7pt]
8(x'''x'''-y'''y''')&=2(x''-y'')^2=&0{,}\mathord{.}\mathord{.}\mathord{.}\mathord{.}\mathord{.}\mathord{.}\mathord{.}\mathord{.}22&3462\\[7pt]
16(x''''x''''-y''''y'''')&=4(x'''-y''')=&0{,}&\\[7pt]
\cline{2-4}
&\mathrm{Summe}\quad=&0{,}0861068379&10
\end{array}
\]
\endgroup
Daraus folgt
\[
\begin{aligned}&\mathrm{d}\mathrm{M}(x,y)=\mathrm{M}(x,y)\times
\Bigl\{\frac{\mathrm{d}x}{2\surd 2}\times1{,}0861068379.\;.\;.\\&\quad+\frac{\mathrm{d}y}{2}\times0{,}9138931621.\;.\;.\Bigr\}\end{aligned}
\]
oder in Zahlen
\[
=\mathrm{d}x\times0{,}460082\mathord{.}\mathord{.}\mathord{.}\mathord{.}
+\mathrm{d}y\times0{,}5474860839
\]
\Needspace{25\baselineskip}
\Needspace{8\baselineskip}\par II. Es sei
\[
\begin{gathered}x=5{,}202778+2{,}784072=7{,}986850,\\[7pt]
y=5{,}202778-2{,}784072=2{,}418706\end{gathered}
\]
wobei \(x,y\) die größten und kleinsten Abstände zwischen Jupiter und Ceres sind (unter Vernachlässigung der Exzentrizität und der Neigung der Ebenen). Daher
\begingroup\small\setlength{\arraycolsep}{3pt}
\[
\begin{array}{lll}
x=7{,}9868500&y=2{,}4187060&xx-y'y'=59{,}\\[7pt]
x'=5{,}2027780&y'=4{,}395207&2(x'x'-y'y')=1{,}28\\[7pt]
x''=4{,}7989925&y''=4{,}781975&4(x''x''-y''y'')=\\[7pt]
x'''=4{,}7904838&y'''=4{,}790476&8(x'''x'''-y'''y''')=\\[7pt]
x''''=4{,}790480&&
\end{array}
\]
\endgroup

\clearpage
\endgroup
